\documentclass[12pt]{article}
\usepackage[a4paper,margin=2.5cm]{geometry}
\usepackage[T1]{fontenc}
\usepackage[utf8]{inputenc}
\usepackage{lmodern}
\usepackage{microtype}
\usepackage{needspace}
\usepackage{booktabs}
\usepackage{amsmath}
\usepackage{xcolor}
\usepackage{listings}
\usepackage[hidelinks]{hyperref}
\definecolor{codegray}{RGB}{245,245,245}
\lstset{basicstyle=\ttfamily\small,backgroundcolor=\color{codegray},frame=single,
  breaklines=true,showstringspaces=false,columns=fullflexible,keepspaces=true}
\setlength{\parindent}{0pt}
\setlength{\parskip}{0.55em}

\begin{document}
\pagenumbering{gobble}
\begin{titlepage}
\centering
{\large Ministerul Educa\c{t}iei \c{s}i Cercet\u{a}rii al Republicii Moldova\par}
\vspace{0.45cm}
{\large Universitatea Tehnic\u{a} a Moldovei\par}
\vspace{0.45cm}
{\large Facultatea Calculatoare, Informatic\u{a} \c{s}i Microelectronic\u{a}\par}
\vspace{0.45cm}
{\large Departamentul Ingineria Software \c{s}i Automatic\u{a}\par}
\vspace{3cm}
{\Large\bfseries Report on Laboratory Work No. 10\par}
\vspace{0.55cm}
{\LARGE\bfseries Add Your Own System Call: trace\par}
\vspace{0.35cm}
{\large Operating Systems\par}
\vfill
\begin{flushleft}
\textbf{Done by:} st. gr. FAF-241, Ca\c{t} Chiril\\[0.8cm]
\textbf{Checked by:} asist. univ. Rei\c{t}man Patricia
\end{flushleft}
\vfill
{Chi\c{s}in\u{a}u -- 2026\par}
\end{titlepage}
\pagenumbering{roman}
\tableofcontents
\clearpage
\pagenumbering{arabic}

\section{Aim, environment, and deliverables}
The aim is to add a system call to xv6 and follow it from the user-side
stub to the kernel handler. In Lab 9, \texttt{strace} observed Linux calls
from outside a program. Here the xv6 kernel itself prints the selected
calls after executing them. Each process stores a bit mask that determines
which syscall numbers are reported.

I implemented all three required parts in the integrated
\texttt{xv6-riscv/} tree. The host tools were Ubuntu 24.04.5 LTS,
RISC-V GCC 13.3.0, and QEMU 8.2.2. The guest was xv6-riscv, with QEMU
configured for three CPUs and 128 MiB of memory. Before editing, I built
and booted the existing tree and ran \texttt{echo baseline-boot-ok}.
The final implementation was rebuilt from clean sources with xv6's
\texttt{-Wall -Werror} compilation settings.

The syscall number is \textbf{23}, because this tree already assigns 22
to \texttt{sync}. Kernel messages use \texttt{printk}, not user-space
\texttt{printf}. The file being searched, \texttt{README}, was left
unchanged at 2441 bytes. Its expected read results are used as an exact
check, not replaced with fabricated output.

The integrated tree is the executable submission. Matching changed source
files and the Makefile are under \texttt{lab10/code/}; the patch records the
changes relative to the starting tree. The report source, assignment,
and recorded runs are under \texttt{report\_files/}. The required
\texttt{notes.txt} preserves the Observe answers. The additional
\texttt{tracetest} program is a regression test, not a second system call.
The optional, ungraded \texttt{sysinfo} extension was not added.

The main build commands on Ubuntu are:
\begin{lstlisting}
sudo apt install build-essential gcc-riscv64-linux-gnu \
  binutils-riscv64-linux-gnu qemu-system-misc bc perl
cd xv6-riscv
make clean
make -j4 TOOLPREFIX=riscv64-linux-gnu- kernel/kernel fs.img
make qemu TOOLPREFIX=riscv64-linux-gnu- CPUS=3
# In QEMU, Ctrl-A followed by X exits the emulator.
\end{lstlisting}
The recorded terminal captures come from real xv6 boots. Process IDs are
specific to those runs; different IDs do not indicate a failure.

\clearpage
\section{Part 1: the number and the user side}
\subsection{The five user-side changes}
\begin{enumerate}
\item Added \texttt{\#define SYS\_trace 23} after \texttt{SYS\_sync} in
\texttt{kernel/syscall.h}.
\item Declared \texttt{int trace(int);} in \texttt{user/user.h}.
\item Added \texttt{entry("trace");} to \texttt{user/usys.pl}. At build
time the script generates the assembly stub in \texttt{user/usys.S}.
\item Created \texttt{user/trace.c} with the handout's argument check,
mask-setting call, and \texttt{exec(argv[2], argv + 2)}.
\item Added \texttt{\$U/\_trace} to the Makefile's \texttt{UPROGS}, so the
program is included in the xv6 disk image.
\end{enumerate}
The generated stub loads \texttt{SYS\_trace} into register \texttt{a7},
executes \texttt{ecall}, and returns. The first argument, the integer mask,
is passed in \texttt{a0}. The program checks that the command and mask are
present, then sets its own mask before replacing its program image:
\begin{lstlisting}[language=C]
if (trace(atoi(argv[1])) < 0) {
  fprintf(2, "trace: failed\n");
  exit(1);
}
exec(argv[2], argv + 2);
fprintf(2, "trace: exec %s failed\n", argv[2]);
exit(1);
\end{lstlisting}
Passing \texttt{argv + 2} makes the requested command its own
\texttt{argv[0]}, with the rest of the command's arguments following it.
Successful \texttt{exec} keeps the same process and does not return to the
failure print.

\subsection{Check 1: reaching an unimplemented syscall}
I built and booted after these five changes, before adding the kernel
handler. The actual output was:
\lstinputlisting{output/checkpoint1.txt}
This checkpoint demonstrates that the executable was on the image, the
stub ran, and the kernel received syscall number 23. The failure was
expected because the dispatch table did not yet contain that handler.
The final submitted sources include the completed kernel side; this
intermediate result is retained separately as evidence.

\clearpage
\section{Part 2: the kernel side}
\subsection{Per-process mask and handler}
I added \texttt{int trace\_mask;} after \texttt{name} in \texttt{struct proc}.
The handler reads argument zero from the saved registers, stores the mask
on the current process, and returns success:
\begin{lstlisting}[language=C]
uint64
sys_trace(void)
{
  int mask;
  argint(0, &mask);
  myproc()->trace_mask = mask;
  return 0;
}
\end{lstlisting}
\texttt{argint} returns \texttt{void} in this version of xv6. The handler
uses the requested integer as a bit mask, not as a syscall number.

\subsection{Dispatch table and post-call printing}
In \texttt{kernel/syscall.c}, I declared \texttt{sys\_trace} and added
\texttt{[SYS\_trace] = sys\_trace} to the handler table. A matching
\texttt{syscallnames} table names every supported call from 1 through 23,
including \texttt{pause}, \texttt{sync}, and \texttt{trace}.
The valid-call path now executes:
\begin{lstlisting}[language=C]
p->trapframe->a0 = syscalls[num]();
if (((uint)p->trace_mask >> num) & 1U)
  printk("%d: syscall %s -> %ld\n", p->pid,
         syscallnames[num], p->trapframe->a0);
\end{lstlisting}
Printing after the handler reports its actual result in \texttt{a0}.
The unsigned shift tests the selected bit; dispatch already checks that the
syscall number is valid. \texttt{\%ld} shows failures as signed \texttt{-1}.
Kernel \texttt{printk} is not a user write syscall and does not trace itself.

\subsection{Initialization, inheritance, and exec}
Allocation initializes the mask to zero; \texttt{freeproc} clears it on
cleanup to avoid stale state in reused slots. This tree's fork implementation
is named \texttt{kfork}. It copies the mask before making the child runnable:
\begin{lstlisting}[language=C]
safestrcpy(np->name, p->name, sizeof(p->name));
np->trace_mask = p->trace_mask;
\end{lstlisting}
The child gets an independent integer mask. \texttt{exec} preserves this
field, so no edit to \texttt{exec.c} is needed. The shell's zero mask still
produces untraced children.

\clearpage
\subsection{Check 2: the required read trace}
\lstinputlisting{output/read.txt}
The observed values exactly match the handout:
\[
1023 + 965 + 453 = 2441\ \text{bytes}.
\]
The final zero indicates end-of-file, not an error. Mask 32 is
\(2^5\), so it selects \texttt{SYS\_read = 5} and suppresses other call
names. In particular, \texttt{open} and \texttt{close} still execute even
though their trace lines are not selected.

The read sizes also agree with the implementation of xv6's \texttt{grep}.
It reads into the unused part of a 1024-byte buffer, reserving one byte;
when no complete line consumes the retained suffix, the next request is
smaller. The three positive returns therefore need not all be 1023.
This run printed no matching README lines, only the selected read traces.

\section{Part 3: playing with the mask}
\subsection{All supported call bits}
\lstinputlisting{output/all.txt}
Mask 2147483647 has bits 0 through 30 set, so it includes all currently
supported syscall numbers. The \texttt{trace} call itself is printed:
its handler enables the mask before the dispatcher checks it. Successful
\texttt{exec} returns 3 here because this implementation returns the new
argument count to become \texttt{main}'s \texttt{argc}; it is not a file
descriptor returned by exec. The subsequent open independently returns
file descriptor 3.

There is no \texttt{exit} line despite its bit being selected. The exit
handler never returns to this post-handler print location. Similarly, the
child's initial return from fork is prepared separately by \texttt{kfork};
the post-dispatch fork line reports the calling parent's return value.
The next subsection demonstrates descendant inheritance through the
children's later calls to fork.

\clearpage
\subsection{Only fork: inherited tracing}
The command \texttt{trace 2 usertests forkforkfork} selects bit 1. The
following is an explicitly shortened excerpt; the complete transcript is
retained in \texttt{output/fork.txt}.
\lstinputlisting[basicstyle=\ttfamily\footnotesize]{output/fork-excerpt.txt}
The trace contained calls from PIDs 5, 6, 7, and 8, not a fixed requirement
of exactly three PIDs. The test runner and its descendants can all call
fork, and copying \texttt{trace\_mask} in \texttt{kfork} lets each caller
print its own selected calls. When allocation fails because the process
table is full, the returned value is \texttt{-1}. The test handles this
condition and finishes with \texttt{ALL TESTS PASSED}. The exact allocation
order and PIDs depend on scheduling.

\subsection{Only write: echo and the newline}
\lstinputlisting{output/write.txt}
\texttt{echo hi} calls \texttt{write} twice: once for the two bytes
\texttt{hi}, then once for the newline. The first write has no newline, so
its trace appears immediately after \texttt{hi}. The second write outputs
a newline before its own trace line. Each trace line also ends with a
newline, which explains the visible blank line. This output ordering is a
consequence of tracing after each write returns, not an extra print from
\texttt{echo}.

\subsection{Selecting open and close together}
The combined mask is
\[
(1 \ll 15)\mathbin{|}(1 \ll 21)
= 32768 + 2097152 = \boxed{2129920}.
\]
\lstinputlisting{output/open-close.txt}
Only the open and close lines were traced; the underlying reads still ran.
The mask is a set of syscall-number bits, not a range of numbers.

\clearpage
\section{Handout questions and answers}
\begin{enumerate}
\Needspace{8\baselineskip}
\item \textbf{Question:} In \texttt{trace 32 grep hello README} why do you see only read, and not open or close?
\par\textbf{Answer:} 32 is \(1 \ll 5\), which enables only the bit for
\texttt{SYS\_read}. Open is 15 and close is 21, so their bits are zero.
Those calls still happen but the dispatcher's mask check does not print them.

\Needspace{8\baselineskip}
\item \textbf{Question:} In forkforkfork, three different pids print lines. Which line of code from Part 2 makes that possible?
\par\textbf{Answer:} \texttt{np->trace\_mask = p->trace\_mask;} in
\texttt{kfork} copies the mask into every child. Descendants therefore trace
their own selected calls. This run showed four caller PIDs (5, 6, 7, and 8),
including the test runner; the inheritance mechanism is the same.

\Needspace{8\baselineskip}
\item \textbf{Question:} In \texttt{trace 65536 echo hi} the trace line is printed between ``hi'' and the new line. Why? (How many times does echo call write?)
\par\textbf{Answer:} Echo calls write twice: 2 bytes for \texttt{hi}, then
1 byte for \texttt{\textbackslash n}. The kernel prints after each completed
write. Consequently the first trace follows \texttt{hi} before echo sends
its newline; the second follows the newline.

\Needspace{8\baselineskip}
\item \textbf{Question:} Which mask traces both open (15) and close (21), and nothing else?
\par\textbf{Answer:} \((1 \ll 15) | (1 \ll 21) = 2129920\).
The recorded combined-mask experiment produced only open and close lines.
The handout's split ``15.'' line is interpreted as syscall number 15 for
open, consistent with \texttt{kernel/syscall.h}.
\end{enumerate}

\section{Additional verification}
I checked a zero mask, an ordinary command after tracing, missing command
arguments, and a nonexistent executable. Mask zero produced no trace lines;
a later untraced command also remained quiet. The missing executable
returned \texttt{-1} from exec. With all bits enabled, xv6's character-wise
\texttt{fprintf} interleaves each error-message character with a write trace;
the verifier removes the trace records before checking the full message.
This was an output-parsing adjustment, not a kernel failure.

\clearpage
\subsection{A focused mask-lifecycle regression}
The additional \texttt{tracetest} enables getpid tracing, forks, and has
the child disable its copy of the mask. The parent waits and calls getpid
again, then disables its own mask. The observed output was:
\lstinputlisting{output/lifecycle.txt}
There are exactly three getpid trace records: the parent before fork, the
child before disabling, and the parent after the child disables. The
untraced getpid calls after each process disables its mask do not add
records. This demonstrates both inheritance and independent per-process
mask changes. The parent also checks the child's exit status.

\subsection{Existing xv6 tests and source consistency}
I ran \texttt{usertests -q}, which executed \textbf{68} tests and finished
successfully. This is xv6's quick suite, not its slow suite. Some memory
protection tests deliberately cause user traps; their diagnostic lines are
expected, and the tests still report success. There were no kernel panics
and no leaked syscall traces in the untraced suite.
\lstinputlisting[basicstyle=\ttfamily\footnotesize]{output/regression-excerpt.txt}
The verifier checks the required traces, masks, error cases, inheritance,
and suite result, and compares all ten submitted source copies with the
tested tree. Raw evidence is under \path{lab10/report_files/output/}.

\section{Conclusion}
Adding trace required both sides of the syscall boundary: a number,
declaration and generated stub in user space, then process state, a handler,
and dispatch-table integration in the kernel. Printing after the handler
makes the returned result visible. Per-process masks persist across exec,
are copied across fork, and can then be changed independently. All required
mask experiments passed, the read byte counts matched the handout exactly,
and the existing xv6 quick suite still passed after the kernel changes.
\end{document}
